\documentclass[10pt]{article}
\providecommand{\FIGDIR}{fig}
\newcommand{\DIGESTDATE}{4 September 2026}
% ===================== PM Research Watch - shared preamble =====================
\usepackage[T1]{fontenc}
\usepackage[utf8]{inputenc}
\usepackage{amsmath,amssymb}
\usepackage{tgpagella}
\usepackage{tgheros}
\renewcommand{\sfdefault}{qhv}
\renewcommand{\ttdefault}{lmtt}
\usepackage{microtype}
\usepackage[english]{babel}

\usepackage[a4paper,top=17mm,bottom=17mm,left=17mm,right=17mm,
            headsep=6mm,footskip=10mm]{geometry}
\usepackage{graphicx}
\usepackage{xcolor}
\usepackage{colortbl}
\usepackage{tikz}
\usetikzlibrary{calc}
\usepackage[most]{tcolorbox}
\usepackage{titlesec}
\usepackage{fancyhdr}
\usepackage{booktabs}
\usepackage{array}
\usepackage{longtable}
\usepackage{enumitem}
\usepackage{multicol}
\usepackage{ragged2e}
\usepackage[hidelinks]{hyperref}

% ------------------------------------------------------------------ palette
\definecolor{Ink}    {HTML}{14262E}
\definecolor{Deep}   {HTML}{0F4C5C}
\definecolor{Teal}   {HTML}{1F7A8C}
\definecolor{Sky}    {HTML}{4FA3C4}
\definecolor{Sage}   {HTML}{6FA287}
\definecolor{Amber}  {HTML}{E3A008}
\definecolor{Clay}   {HTML}{B5643C}
\definecolor{Coral}  {HTML}{D1495B}
\definecolor{Violet} {HTML}{6A5B9E}
\definecolor{Slate}  {HTML}{4A6572}
\definecolor{Paper}  {HTML}{FBF9F5}
\definecolor{Mist}   {HTML}{EEEAE2}
\definecolor{Rule}   {HTML}{D9D3C9}

\pagecolor{Paper}
\color{Ink}

\hypersetup{colorlinks=true, linkcolor=Deep, urlcolor=Teal, citecolor=Teal,
            pdftitle={PM Research Watch}, pdfauthor={Automated daily digest}}

% ------------------------------------------------------------------ headings
\titleformat{\section}
  {\sffamily\bfseries\large\color{Deep}}{}{0pt}
  {}[{\vspace{1.2mm}\color{Amber}\titlerule[1.1pt]}]
\titlespacing*{\section}{0pt}{5.5mm}{2.8mm}

\titleformat{\subsection}
  {\sffamily\bfseries\normalsize\color{Teal}}{}{0pt}{}
\titlespacing*{\subsection}{0pt}{3.6mm}{1.4mm}

\setlist[itemize]{leftmargin=4.2mm, itemsep=1.1mm, topsep=1.3mm, parsep=0pt,
                  label=\textcolor{Teal}{\footnotesize$\blacktriangleright$}}

% ------------------------------------------------------------------ header/footer
\pagestyle{fancy}
\fancyhf{}
\renewcommand{\headrulewidth}{0pt}
\renewcommand{\footrulewidth}{0pt}
\fancyhead[L]{\sffamily\scriptsize\color{Slate}PM RESEARCH WATCH\;\textcolor{Rule}{\textbar}\;\DIGESTDATE}
\fancyhead[R]{\sffamily\scriptsize\color{Slate}Particulate matter monitoring \& health effects}
\fancyfoot[C]{\sffamily\scriptsize\color{Slate}\thepage}

% ------------------------------------------------------------------ components
% metric tile
\newcommand{\metric}[3]{%
  \begin{tcolorbox}[enhanced, colback=white, colframe=Rule, boxrule=0.5pt,
    arc=1.6mm, left=2.2mm, right=2.2mm, top=1.8mm, bottom=1.8mm,
    borderline west={1.7mm}{0pt}{#3}, nobeforeafter, valign=center,
    height=17mm, width=\linewidth]
    \RaggedRight
    {\sffamily\bfseries\fontsize{17}{17}\selectfont\color{#3}#1}\\[0.5mm]
    {\sffamily\fontsize{6.5}{7.4}\selectfont\color{Slate}#2}
  \end{tcolorbox}}

% pull-quote / key insight box
\newtcolorbox{keybox}[1][]{enhanced, colback=Mist, colframe=Mist,
  arc=1.6mm, left=3.2mm, right=3.2mm, top=2.4mm, bottom=2.4mm,
  borderline west={1.7mm}{0pt}{Amber}, #1}

% subtopic band
\newcommand{\band}[2]{%
  \vspace{3.2mm}\noindent
  \begin{tcolorbox}[enhanced, colback=#2, colframe=#2, arc=1.2mm,
    left=2.6mm, right=2.6mm, top=1.1mm, bottom=1.1mm, boxrule=0pt,
    width=\linewidth, nobeforeafter]
    {\sffamily\bfseries\small\color{white}#1}
  \end{tcolorbox}\vspace{1.6mm}}

% one paper entry: {short cite}{title}{journal}{doi}{body}
\newcommand{\paperentry}[5]{%
  \par\vspace{1.9mm}\noindent
  {\sffamily\bfseries\footnotesize\color{Ink}#1}\;%
  {\sffamily\fontsize{7}{8}\selectfont\color{Slate}#3}\\[0.4mm]
  {\footnotesize\itshape\color{Deep}#2}\\[0.6mm]
  {\footnotesize #4}\\[0.3mm]
  {\fontsize{6.6}{7.6}\selectfont\color{Slate}\texttt{doi:}\,\href{https://doi.org/#5}{\texttt{\detokenize{#5}}}}
  \par}

% caption style
\newcommand{\figcap}[1]{\par\vspace{1.0mm}
  {\fontsize{7.2}{8.6}\selectfont\color{Slate}\textsf{\textbf{Fig.}}\; #1\par}\vspace{1.5mm}}

\providecommand{\permil}{\textperthousand}

\setlength{\parindent}{0pt}
\setlength{\parskip}{0pt}
\renewcommand{\arraystretch}{1.18}


\begin{document}

% ==================================================================== MASTHEAD
\thispagestyle{empty}
\begin{tikzpicture}[remember picture, overlay]
  \fill[Deep] (current page.north west) rectangle
      ($(current page.north east)+(0,-67mm)$);
  \fill[Amber] ($(current page.north west)+(0,-67mm)$) rectangle
      ($(current page.north east)+(0,-68.4mm)$);
  \foreach \x/\y/\r in {22/12/0.9, 41/26/1.5, 63/9/0.7, 88/31/1.1, 112/17/1.9,
                        138/29/0.8, 159/13/1.3, 178/54/1.0, 196/21/1.6,
                        30/44/0.6, 74/57/1.2, 125/48/0.9, 168/6/0.8, 191/59/0.7,
                        52/18/1.0, 100/60/0.6, 148/20/1.1, 66/49/1.4, 205/45/0.8}
    \fill[Sky, opacity=0.30] ($(current page.north west)+(\x mm,-\y mm)$) circle (\r mm);
\end{tikzpicture}

\vspace*{4mm}
{\sffamily\fontsize{7.6}{9}\selectfont\color{Sky}%
 AUTOMATED DAILY LITERATURE DIGEST\;\textbullet\;PhD RESEARCH WATCH\par}
\vspace{2.2mm}
{\sffamily\bfseries\fontsize{27}{29}\selectfont\color{white}%
 Particulate Matter\par}
{\sffamily\fontsize{27}{29}\selectfont\color{Sky}%
 Monitoring \& Health Effects\par}
\vspace{5.0mm}
{\sffamily\fontsize{8.4}{10}\selectfont\color{white}%
 Issue 2026-09-04\;\;\textcolor{Amber}{\textbar}\;\;Entry window 4 September 2026 (single day)\;\;%
 \textcolor{Amber}{\textbar}\;\;27 records in scope, 55 screened out\par}
\vspace{18mm}

% -------------------------------------------------------------- metric strip
\noindent
\begin{minipage}[t]{0.190\linewidth}\metric{27}{RECORDS IN SCOPE\\(55 SCREENED OUT)}{Deep}\end{minipage}\hfill
\begin{minipage}[t]{0.190\linewidth}\metric{9}{SUBTOPIC CLUSTERS\\OF TEN OCCUPIED}{Teal}\end{minipage}\hfill
\begin{minipage}[t]{0.190\linewidth}\metric{11}{EFFECT ESTIMATES\\(FIVE SOURCES)}{Amber}\end{minipage}\hfill
\begin{minipage}[t]{0.190\linewidth}\metric{6}{SENSING /\\INSTRUMENTATION}{Clay}\end{minipage}\hfill
\begin{minipage}[t]{0.190\linewidth}\metric{12}{RECORDS WITH NO\\HEALTH ENDPOINT}{Sky}\end{minipage}

\vspace{6mm}

% -------------------------------------------------------------- signal of day
\section*{\sffamily\bfseries\color{Deep}Signal of the day}
\vspace{-1.5mm}{\color{Amber}\rule{\linewidth}{1.1pt}}\vspace{2.5mm}

\begin{keybox}
{\sffamily\bfseries\small\color{Deep}Two papers, two countries, one result: the two
quantities a PM instrument actually measures --- scattered light and gravimetric mass ---
both decouple from the fraction that does the damage, and both fail in the direction that
under-reports it.}\\[1.4mm]
\footnotesize
Zeng et al.\ (\textit{J.\ Hazard.\ Mater.}) put volume-normalised oxidative potential and
light extinction on a common quadrant plane at two urban sites in southern Taiwan and
isolate an \textbf{``invisible killer'' regime: high redox activity, low optical impact}.
The mediator is thermodynamic, not chemical --- \textbf{high aerosol liquid water content}
drives hygroscopic growth into optics-dominant states, while \textbf{low ALWC} corresponds
to the highest OP-to-extinction ratios; secondary nitrate loads the extinction side
specifically. Independently, Zhao C et al.\ (\textit{Environ.\ Pollut.}) tracked 16 priority
particle-bound PAHs across a real air-quality improvement in Kunming. Mass fell
\textbf{32.7\%} and BaP-equivalent burden \textbf{31.4\%}, but the
\textbf{particle-normalised toxicity, BaP$_{\mathrm{eq}}$/PM$_{2.5}$, did not move
($-$7.6\%, $p$=0.274)}; 4--6-ring species stayed above \textbf{81\%} of $\Sigma$PAH, and
\textbf{40.3\%} of simulated adult lifetime cancer risks still exceeded $10^{-6}$.
Put together: a nephelometric sensor reports the wrong thing when the air is dry, a
gravimetric filter reports the wrong thing when the air is cleaner, and the error is the
same sign in both cases. A network built to certify compliance against a mass threshold
will read a genuine improvement while the per-microgram hazard is unchanged. Neither paper
is a sensor study, and that is the point --- this is the calibration target problem stated
from the chemistry side.
\end{keybox}

\vspace{3mm}

% -------------------------------------------------------------- exec brief
\section*{\sffamily\bfseries\color{Deep}What changed today}
\vspace{-1.5mm}{\color{Amber}\rule{\linewidth}{1.1pt}}\vspace{2.5mm}

\footnotesize
\begin{itemize}
\item \textbf{A normal-sized, unusually well-formed issue after two outliers.} 27 records
against 38 on 3 September and 33 on 1 September. \textbf{24 distinct design labels over 27
records}, and for the first time in this run of issues \textbf{not one required a new entry
in the design map} --- 3 September needed fourteen. The corpus is broad without being
ragged: \textit{Measurement campaign} (7) leads, then \textit{Review / synthesis} and
\textit{Observational --- cohort} at 4 each.

\item \textbf{The no-endpoint share fell hard, and that is a composition change rather than
a trend.} \textbf{12 of 27 (44\%)} carry no human health endpoint, against 61\%, 60\% and
58\% on the three preceding issues. The sensing block halved (11 $\rightarrow$ 6) while
cohort epidemiology held, so the ratio moved because the \textit{ES\&T} by-ISSN deposit
that inflated 3 September did not repeat.

\item \textbf{Two UK Biobank cardiovascular papers, and their intervals are not on the same
evidential footing.} Chen N et al.\ report per-IQR incident-VTE hazard ratios of
\textbf{1.052 (1.005--1.101)} for elemental carbon and \textbf{1.090 (1.010--1.176)} for
ammonium on 9{,}988 events --- both lower bounds sit within 0.01 of the null, which is the
honest width for a time-invariant residential exposure. The same paper's heatwave hazard
ratios run \textbf{1.149 (1.137--1.161)} to \textbf{1.407 (1.396--1.417)}: intervals
roughly \textbf{three times narrower} on the same events, because a daily time-varying
exposure carries far more information than an annual surface. Both belong in the paper;
they should not be read off one forest panel as comparable quantities, and the mediation
claim (elemental carbon and ammonium partly mediating the greenness effect) inherits the
weaker exposure.

\item \textbf{Two prenatal-PM cohorts with exactly inverted strengths.} Mbishi et al.\ have
\textbf{119{,}110} Flemish children and municipality-scale kriged PM, and report a
head-circumference deficit at 12 months of \textbf{$-$0.065 $z$ per IQR PM$_{2.5}$} --- about
one-fifteenth of a standard deviation, unambiguous at that $n$ and, as the authors
themselves write, small. Wei Q et al.\ have \textbf{2{,}680} pairs and weekly
component-resolved exposure, and can name windows the first study cannot see: femur length
most sensitive to black carbon at \textbf{gestational weeks 12--15}, head circumference to
nitrate at \textbf{weeks 5--12 and 22--24}. Power and exposure resolution are being traded
against each other, and no record today has both.

\item \textbf{The widest source spread on the forest panel since August.} 11 estimates from
\textbf{five} sources, tying 27 August; only the 26--27 July back-scan and 6 August have had
more. One source still supplies 6 of the 11, so the panel is more diverse than it is
balanced.

\item \textbf{Geography narrowed.} China 8 of 27 (30\%), \textit{Global / multi-region} 7
(of which four are genuine multi-region reviews, one is a laboratory record and two are
metadata-only records with no stated setting), East Asia ex-China 4. \textbf{Sub-Saharan
Africa, South Asia, the Middle East and Oceania are all absent}, three days after
Sub-Saharan Africa tied its watch record at three. Latin America returns with a single
Merida aerosol-chemistry record.

\item \textbf{One new interventional trial registration, and it is a registry-hygiene
warning.} NCT07805824 is a cluster RCT of a workplace heat and air-pollution mitigation
package in \textbf{20 restaurants / 200 oven workers} in Najaf, Iraq, with
\textbf{personal PM$_{2.5}$ by gravimetric and real-time monitoring} as its sole secondary
endpoint --- directly relevant to occupational personal sensing. But it is filed as
\textbf{Phase 4} with no drug, its start and primary-completion dates are both 2026-01-01
while its status is already \textsc{completed}, and the public record still carries the
submitter's own pre-submission checklist asking them to ``confirm the exact first-participant
enrollment date''. Treat the endpoint design as usable and the timeline fields as not.

\item \textbf{Nil legs, one of them now verified rather than inspected.} Consensus is nil
for the \textbf{sixth} consecutive run, but this time by DOI-level date check: its three
low-cost-sensor calibration hits are all real and all out of window (Crossref
\texttt{created} 2026-03-26 and 2026-04-21), and the third --- the Alphasense/Sensirion
recalibration-cadence paper --- \emph{is already in this corpus} from an earlier issue. The
leg finds the right literature; only the date filter is missing.
\textbf{OpenAlex} HTTP 429, unchanged since 28 July. \textbf{arXiv} answered cleanly for the
second consecutive run, newest submission still \textbf{2026-08-20}, contribution again zero.
\end{itemize}

\vspace{3mm}

% ==================================================================== ANALYTICS
\section*{\sffamily\bfseries\color{Deep}Corpus analytics}
\vspace{-1.5mm}{\color{Amber}\rule{\linewidth}{1.1pt}}\vspace{3mm}

\noindent\includegraphics[width=\linewidth]{\FIGDIR/f1_subtopics.png}
\figcap{Subtopic assignment of the 27 in-scope records, single-label by dominant
contribution. \textbf{Nine of ten clusters are occupied}, which has now happened on 12 of
41 issues; the vacant one is \textit{Burden, policy \& mitigation}, and it is vacant
because today's two burden-adjacent records (US building-sector mortality, transboundary
smoke exposure) both lead with an exposure surface and were filed under
\textit{Exposure assessment \& modelling} instead. \textit{Sensing} (6) and
\textit{Exposure} (4) hold 10 of 27, against 23 of 38 yesterday --- the instrument-heavy
run has broken.}

\vspace{2mm}
\noindent\includegraphics[width=\linewidth]{\FIGDIR/f2_design_endpoint.png}
\figcap{Left --- study architecture over \textbf{24 distinct design labels}, every one of
which was already in \texttt{plots.py}'s design map; no map edit was needed this run,
against fourteen additions on 3 September. \textit{Measurement campaign} (7) leads, ahead
of \textit{Review / synthesis} and \textit{Observational --- cohort} at 4 each. Only
\textbf{2} records ship as \textit{Metadata only}, both Elsevier titles whose publishers
deposited no abstract anywhere the pipeline can reach. Right --- \textbf{12 of 27 (44\%)}
carry no human health endpoint, the lowest share since 29 August and well below the 61\%
of the preceding issue. Behind that bar the health endpoints are unusually flat: none
exceeds 4, and \textit{Cardiovascular} reaches 4 only by pooling two cohorts, a time-series
and a green-space cross-section. \textit{Cancer} (1) and \textit{Oncologic} (1) are one
quantity drawn as two bars --- an oesophageal incidence rate ratio and a computed
carcinogenic-risk index --- because \texttt{ENDPOINT\_CANON} still does not merge them.
That defect was logged on 3 September and is unfixed.}

\vspace{2mm}
\noindent\includegraphics[width=\linewidth]{\FIGDIR/f3_metric_geography.png}
\figcap{Left --- particle metric. \textit{PM (unspeciated) / emissions} at 9 carries the
records whose measurand is not a regulated mass fraction: BVOC fluxes, odour activity
values, bioaerosol counts, surface-attachment fractions and the two abstract-free records.
\textit{PM$_{2.5}$ only} (6) plus \textit{PM$_{2.5}$ + PM$_{10}$ jointly} (4) is the
epidemiological core; \textit{Composition / speciation} (3) is where today's signal sits.
Right --- geography. China 8, \textit{Global / multi-region} 7, East Asia ex-China 4
(two South Korean, one Japanese, one Taiwanese), Europe 3, North America 2, Southeast
Asia 2, Latin America 1. \textbf{Sub-Saharan Africa, South Asia, the Middle East and
Oceania are all empty.} North America at 2 continues a thin run --- 1 on 27 August, 2 on
1 September, 8 on 3 September --- so the 3 September figure was the anomaly, not this one.}

\vspace{2mm}
\noindent\includegraphics[width=\linewidth]{\FIGDIR/f6_lifecourse.png}
\figcap{Life-course windows addressed; counts are non-exclusive and cover only records with
a human endpoint or an exposed human population. \textbf{Every stage is occupied}, which is
the first time since 27 August --- puberty and adolescence, empty on the two preceding
issues, is filled by a single record (18{,}398 Chinese adolescents). Working-age adulthood
(4) leads, carried by the two UK Biobank cohorts, the welding-fume review and the Ontario
household survey. ``Older adults'' is generous here: the two records are age
$\geq$\,60 (Thailand) and age $\geq$\,55 (GBD), not 65+.}

\vspace{2mm}
\noindent\includegraphics[width=\linewidth]{\FIGDIR/f4_heatmap.png}
\figcap{Subtopic $\times$ study-architecture cross-tabulation. The \textit{Sensing} row is
entirely \textit{Measurement campaign} and \textit{Chamber / laboratory} today --- no
modelling or forecasting records were filed under sensing, which is why the cluster is
smaller but purer than yesterday's. \textit{Cardiovascular \& metabolic} and
\textit{Other clinical endpoints} are the only subtopics spread across three architecture
families, and one of the three in the cardiovascular row is \textit{Metadata only}, so part
of that apparent breadth is an artefact of a missing abstract rather than real variety.}

\vspace{2mm}
\noindent\includegraphics[width=\linewidth]{\FIGDIR/f5_forest.png}
\figcap{Eleven estimates from \textbf{five} sources --- the widest source spread since
27 August. They are not commensurable and should not be read as a pooled panel. Six are
from one UK Biobank VTE analysis and are internally correlated: three per-IQR PM$_{2.5}$
component hazard ratios whose lower bounds sit within 0.01 of the null, one protective
greenness hazard ratio, and two heatwave hazard ratios whose intervals are roughly three
times narrower because the exposure is daily rather than annual. The two ozone rate ratios
for out-of-hospital cardiac arrest are identical point estimates in men and women (1.22)
with different widths, which is what sex-stratification of 7{,}996 events buys and nothing
more. The single \textbf{IRR 1.055 (0.995--1.119)} for respiratory visits crosses the null
and is reported as null by its own authors. The black-carbon self-injury odds ratio
\textbf{1.76 (1.60--1.94)} is the largest effect on the panel and the weakest design ---
one cross-sectional survey round, exposure assigned by residence. Beta coefficients from
the two fetal-growth cohorts are absent: a $z$-score mean difference has no place on a
log-ratio axis.}

\clearpage

% ==================================================================== DIGEST
\section*{\sffamily\bfseries\color{Deep}Paper-level digest}
\vspace{-1.5mm}{\color{Amber}\rule{\linewidth}{1.1pt}}\vspace{1mm}

{\fontsize{7.4}{9}\selectfont\color{Slate}Ordered by cluster. Each entry gives the
methodological core and the finding that matters, not an abstract paraphrase.\par}

\band{Sensing, forecasting \& instrumentation\hfill 6 records}{Amber}

\paperentry{Zeng et al.\ 2026d}
{Oxidative potential and light extinction decouple, and aerosol water is the switch}
{J Hazard Mater}
{Paired oxidative-potential and optical measurements at two urban sites in southern
Taiwan, arranged into a quadrant framework on \emph{volume-normalised} OP against light
extinction. Under coupling, mass-normalised OP and mass extinction efficiency share their
drivers, principally \textbf{water-soluble organic carbon}. Under decoupling, \textbf{high
aerosol liquid water content} pushes the aerosol into an optics-dominant state through
hygroscopic growth, while \textbf{low ALWC} yields the highest OP-to-extinction ratios ---
the ``invisible killer'' quadrant. Secondary nitrate loads extinction specifically; aerosol
pH is a weaker secondary term. The direct consequence for distributed sensing: every
nephelometric low-cost sensor estimates mass from scattering, so its error against the
redox-active fraction is not random but conditioned on humidity, in the same direction as
its already-known RH bias. Weakness: OP is one assay on one aerosol population at two
sites, sampling duration is not stated in the abstract, and no causal chain from OP to a
clinical endpoint is tested here.}
{10.1016/j.jhazmat.2026.143459}

\paperentry{Ichikawa et al.\ 2026}
{Drone sampling agrees with a flux tower to within 30\%, and that is the useful number}
{Atmos Chem Phys}
{Multi-year warm-season BVOC measurements at multiple heights on a 30\,m flux tower in a
\textit{Quercus serrata}-dominated suburban Tokyo forest, with supplemental drone sampling
and comparison against MEGAN. Isoprene is \textbf{over 90\% of measured BVOC composition}
at peak summer, monoterpenes stay low with weak vertical gradients, and daily average
isoprene emission fluxes span \textbf{$-$0.05 to 15.30\,mg\,m$^{-2}$\,h$^{-1}$}. The
methodological result is the one worth carrying: \textbf{drone-measured horizontal
variability is 10--30\% within 30\,m of the tower}, and \textbf{tower- and drone-derived
flux estimates differ by about 30\%}. That is a direct, published estimate of the
representativeness error a mobile aerial platform introduces relative to a fixed mast ---
the same quantity a mobile PM platform needs and rarely reports. Limits: BVOCs are
SOA precursors, not particles; intermittent campaigns over multiple years confound
season with instrument configuration.}
{10.5194/acp-26-12521-2026}

\paperentry{Kim C-h et al.\ 2026}
{Block-missing monitoring data, five imputation models, and an honest sensitivity test}
{Atmosphere}
{Eleven months of 20-minute measurements of 22 designated odour compounds at a livestock
farm, a wastewater plant and an organic-waste plant in Eumseong, South Korea, with
\textbf{multi-week to multi-month block gaps} --- the failure mode that actually afflicts
long-running low-cost networks, as opposed to the scattered single-point dropout most
imputation papers assume. Compound-specific performance was assessed by \textbf{block
holdout} against BiLSTM, BiGRU, TCN, Transformer and HistGradientBoosting. Odour activity
values were computed \emph{primarily from observed concentrations}, with the fully imputed
series used only as a sensitivity comparison to quantify interpolation-induced bias, which
is the correct discipline and is rarely followed. Ammonia was excluded outright for a low
valid-observation rate. Odour contribution was dominated by volatile fatty acids and
trimethylamine rather than hydrogen sulfide. Transferable to PM networks in method, not in
analyte.}
{10.3390/atmos17090868}

\paperentry{Cai J et al.\ 2026}
{Subway air microbiome is homogeneous at network scale and heterogeneous at platform scale}
{Environ Sci Technol}
{Stratified cluster sampling of \textbf{43 underground stations across 13 lines} of the
Shanghai subway, collecting total suspended particles and PM$_{2.5}$ \emph{simultaneously}
at entrances, halls and platforms, then integrating taxonomy, transcriptional potential,
PM$_{2.5}$ chemistry, microclimate and network topology. Bacteria dominate both size
fractions and most transcriptional signal; eukaryotes are strongly enriched in TSP,
\textbf{a clean demonstration of particle-size filtering of the biological load}. Dominant
taxa (\textit{Staphylococcus}, \textit{Bacillus}) show weak abundance--transcription
coupling, so abundance is a poor proxy for activity. Diversity and composition are
essentially uniform across the network --- consistent with passenger mixing --- while
finer-scale structure appears within stations. For exposure work this means a single
station is a reasonable network sentinel for community composition and a poor one for
microenvironment. Weakness: one city, one season not stated, no human outcome.}
{10.1021/acs.est.6c00174}

\paperentry{Rosas et al.\ 2026}
{The PM$_{2.5}$/PM$_{10}$ ratio separates two aerosol regimes that mass alone cannot}
{Environ Sci Pollut Res Int}
{Continuous PM$_{2.5}$, PM$_{10}$ and particle-bound PAH at a near-coastal urban site in
Merida, Mexico, joined to FIRMS fire detections and HYSPLIT back-trajectories, contrasting
the March--May regional biomass-burning season against the June--August African dust
season. Concentrations are the wrong discriminator: \textbf{PM$_{2.5}$ and PM$_{10}$ rise
only modestly during dust}, but the \textbf{PM$_{2.5}$/PM$_{10}$ ratio falls}, marking the
coarse mineral contribution, while biomass burning loads the fine mode and the combustion
tracers. A network reporting PM$_{2.5}$ mass alone would score the dust season as the
milder of the two and miss the toxicologically distinct composition entirely. Part I of a
two-part study, so the health framing in the title is anticipatory rather than delivered
here. Single site, single urban domain.}
{10.1007/s11356-026-38193-x}

\paperentry{Seong \& Hoque 2026}
{Attachment fraction ranks metal $>$ wood $>$ glass, and carpet is two microenvironments}
{Front Pediatr}
{Batch attachment and centrifugal detachment of \textit{Corynebacterium} sp.\ on carpet,
wood, metal and glass, with near-surface velocity fields resolved by particle image
velocimetry and surfaces characterised for roughness, contact angle and porosity.
Attachment fraction rose with exposure time on all materials, reaching \textbf{$\sim$0.59
on metal, $\sim$0.48 on wood, $\sim$0.36 on glass}. The useful observation is
morphological: \textbf{carpet generates two distinct microenvironments} --- a
higher-velocity outer-fibre region and a sheltered low-velocity inter-fibre region --- so a
single resuspension coefficient for ``carpet'' is a category error. Wood produced the most
spatially heterogeneous velocity field, glass the most uniform. Relevance to PM is
indirect but real: resuspension source terms in indoor exposure models are usually fitted
per material with no flow-field justification. One organism, no PM measurement, replicate
count not stated.}
{10.3389/fped.2026.1895305}

\band{Exposure assessment \& modelling\hfill 4 records}{Amber}

\paperentry{Ullah et al.\ 2026}
{Transboundary smoke moves almost nobody across China's standard and almost everybody
across WHO's}
{Environ Pollut}
{Daily PM$_{2.5}$ from \textbf{51 cities} in Yunnan, Guangxi and Guangdong for
March--April 2015--2019, gridded onto WorldPop, with the Indo-China Peninsula fire
contribution isolated by paired GEOS-Chem runs with and without the fires. The
population-weighted share of days above the \textbf{WHO 2021 24\,h guideline
(15\,$\mu$g/m$^3$) reached 78--96\% every year}; above \textbf{China's Grade II standard
(75\,$\mu$g/m$^3$) it was 1--6\%}, an average daily exposed population of about
\textbf{6 million}, under 3\% of the corridor's residents. The two benchmarks describe
different worlds from identical data, which is the paper's real contribution to monitoring
policy. Principal weakness: the attribution runs at \textbf{2$^\circ$ $\times$
2.5$^\circ$}, roughly 200\,km, so the fire contribution is a regional average imposed on
city-scale observations and cannot support any within-province ranking.}
{10.1016/j.envpol.2026.129081}

\paperentry{Arunachalam et al.\ 2026}
{Building-sector combustion mortality is three times prior estimates, and it is mostly NO$_2$}
{Research Square (preprint)}
{CMAQ Decoupled Direct Method sensitivities of NO$_2$, O$_3$ and PM$_{2.5}$ to
commercial and residential building-sector precursor emissions for the contiguous US
(2021), propagated through BenMAPR. Net burden \textbf{$\approx$14{,}000 premature deaths
per year (95\% CI 7{,}900--19{,}000)}: \textbf{12{,}000 from NO$_2$}, \textbf{3{,}100 from
PM$_{2.5}$}, and \textbf{1{,}000 deaths avoided} through O$_3$ reductions. Residential gas
alone accounts for \textbf{$\approx$7{,}900}. The headline ``three times higher than
previously reported'' is carried almost entirely by the NO$_2$ term, so the result is a
statement about which NO$_2$ concentration--response function you accept, not about
building emissions being larger than thought. Preprint, not peer reviewed. The negative
O$_3$ term is the sign of a titration effect and is correctly reported rather than
suppressed, which is a point in the paper's favour.}
{10.21203/rs.3.rs-10758198/v1}

\paperentry{Zhao C et al.\ 2026}
{Mass fell a third; toxicity per microgram did not move}
{Environ Pollut}
{Sixteen priority PM$_{2.5}$-bound PAHs in Kunming, compared between a higher- and a
lower-PM$_{2.5}$ stage. Raw declines: \textbf{PM$_{2.5}$ $-$32.7\%,
$\Sigma$PAHs $-$27.4\%, BaP$_{\mathrm{eq}}$ $-$31.4\%}; covariate-adjusted,
\textbf{$-$28.8\%} and \textbf{$-$33.5\%}. But the particle-normalised intensity
\textbf{BaP$_{\mathrm{eq}}$/PM$_{2.5}$ changed by $-$7.6\% ($p$=0.274)} and
BaP$_{\mathrm{eq}}$/$\Sigma$PAH by only $-$6.5\% ($p$=0.026). Four-to-six-ring species held
above \textbf{81\%} of $\Sigma$PAH and BaP plus dibenzo[a,h]anthracene supplied over
\textbf{82\%} of BaP$_{\mathrm{eq}}$ in both stages; \textbf{40.3\%} of simulated adult
incremental lifetime cancer risks still exceeded $10^{-6}$. PMF patterns are read
conservatively as source-\emph{related} rather than as resolved sources, which is the
right caution. The inference is threatened by the two-stage design: ``stage'' bundles
season, meteorology and policy, and no counterfactual separates them.}
{10.1016/j.envpol.2026.129094}

\paperentry{Leung et al.\ 2026}
{Abandoned wells, indoor VOCs, and a pilot honest enough to report its own limits}
{Indoor Air}
{Passive indoor sampling for 13 VOCs in \textbf{17 households (October 2024)} and
\textbf{15 households (February--March 2025)} in southwestern Ontario, with abandoned-well
exposure estimated by inverse-distance-squared weighting over wells within 3\,km and
self-reported acute symptoms over the previous six months and during sampling. This is
gas-phase, not particulate --- carried here because the exposure-assignment design is the
one a distributed indoor PM study would use and because the MAPSH project measures both.
At $n$$<$20 per round with Spearman correlations and no personal monitoring, the study is
powered to describe indoor concentrations against Canadian guidelines and not to detect a
well-proximity association; the six-month symptom recall window compounds that. Read as a
feasibility and concentration-range report.}
{10.1155/ina/8041546}

\band{Occupational \& indoor\hfill 4 records}{Clay}

\paperentry{Ibrahim 2026b}
{The case for routine indoor PM monitoring, argued without data}
{Qeios (preprint)}
{A commentary arguing that Malaysia's ICOP IAQ 2010 framework, though it contains
preventive elements, is operated as a complaint-triggered investigation regime, and that
\textbf{particulate matter and outdoor-to-indoor pollutant transport} should join CO$_2$ as
routinely monitored parameters in building management. The specific climatic argument ---
that hot, humid conditions couple ventilation, condensation and biological growth in ways a
temperate-derived code handles badly --- is sound and underexploited in the LCS literature.
Everything else is assertion: single author, non-peer-reviewed platform, no primary data,
and the Malaysian office, hospital and school studies it cites are described rather than
synthesised. Included as an argument about monitoring design, not as evidence.}
{10.32388/b506md}

\paperentry{Wei J et al.\ 2026}
{Welding fume immunology, with the exposure-assessment gap named explicitly}
{Front Immunol}
{Mini review of metal-rich fine and ultrafine welding-fume particles as a stimulus to
airway epithelium, alveolar macrophages, neutrophils, monocytes and circulating immune
cells. Consistent signals across human exposure studies, controlled inhalation, cell models
and workplace investigations: acute-phase proteins, cytokines, chemokines and neutrophilic
responses, with some evidence of altered T-helper activity and blunted responsiveness to
bacterial stimuli. TRPA1/TRPV1 sensory signalling is proposed as a rapid route to cough and
airway reflex, and the review is careful to say that direct welding-fume-specific human
evidence for that pathway is limited. The stated obstacle to synthesis is the one that
matters here: \textbf{fume composition and particle metric vary so much between studies
that biomarker results cannot be pooled}. No quantitative estimate, no pooling, no
systematic search protocol reported.}
{10.3389/fimmu.2026.1882845}

\paperentry{Kusi et al.\ 2026}
{Waste-tyre combustion toxicology: a systematic search that could not pool}
{Research Square (preprint)}
{PRISMA-style search of PubMed, Scopus, ScienceDirect and Google Scholar for 2016--2025
human and experimental studies of waste-tyre combustion emissions --- particulate matter,
PAHs, VOCs and heavy metals --- reporting biochemical, molecular, biomonitoring or
toxicological outcomes, appraised with SYRCLE, JBI and NIH instruments as appropriate to
design. \textbf{Heterogeneity in exposure conditions, designs and reported outcomes
precluded quantitative synthesis}, which the authors state plainly. The value is in
mapping organ-specific mechanistic pathways for an exposure that is ubiquitous in informal
waste economies and almost absent from regulated monitoring networks. Preprint, not peer
reviewed; a narrative synthesis of a heterogeneous literature inherits every selection
effect in that literature.}
{10.21203/rs.3.rs-10530000/v1}

\paperentry{Derwein et al.\ 2026}
{Masks, filtration and ventilation on exhaled-air transmission --- metadata only}
{Build Environ}
{No abstract was deposited to Crossref, Europe PMC or Semantic Scholar at index time, and
ScienceDirect refused retrieval from this environment, so nothing beyond title, authorship
and venue can be asserted. Recorded so the DOI enters \texttt{seen.json} and the record
is not re-screened; queued for the weekly abstract re-check. Listed under
\textit{Occupational \& indoor} on the title's face, which is a provisional assignment.}
{10.1016/j.buildenv.2026.115206}

\band{Cardiovascular \& metabolic\hfill 4 records}{Teal}

\paperentry{Chen N et al.\ 2026}
{Elemental carbon leads the PM$_{2.5}$ components on venous thromboembolism, barely}
{Eur J Prev Cardiol}
{\textbf{483{,}610} UK Biobank participants (mean age 56.5, 54.3\% female), median
\textbf{11.83} years, \textbf{9{,}988} incident VTE from hospital and death records.
Per-IQR hazard ratios: elemental carbon \textbf{1.052 (1.005--1.101)}, sulfate
\textbf{1.064 (1.001--1.130)}, ammonium \textbf{1.090 (1.010--1.176)}; organic matter and
nitrate not significant. Residential greenness \textbf{0.944 (0.916--0.972)}, partly
mediated by elemental carbon and ammonium. Heatwave hazard ratios span \textbf{1.149
(1.137--1.161)} to \textbf{1.407 (1.396--1.417)} across definitions. Two cautions. The
three positive component intervals all have lower bounds within 0.01 of unity at
3$\times$3\,km resolution, so component \emph{ordering} is not established. And the
heatwave intervals are roughly three times narrower on the same events because that
exposure is daily rather than annual --- a real difference in information, not a real
difference in certainty about air pollution.}
{10.1093/eurjpc/zwag465}

\paperentry{Fang et al.\ 2026}
{Lifestyle matters more where the air is worse, and most for early-onset heart failure}
{Eur J Prev Cardiol}
{\textbf{233{,}341} UK Biobank participants, 2006--2010 baseline to 2023, with an air
pollution score over NO$_2$, NO$_x$, PM$_{10}$, PM$_{2.5}$ and PM$_{2.5-10}$ crossed
against a six-component lifestyle score. Additive interaction is present and larger for
early onset: RERI \textbf{0.44 (0.14--0.74)} for all heart failure, \textbf{1.33
(0.03--2.62)} early-onset, \textbf{0.39 (0.08--0.70)} late-onset, with the early-versus-late
contrast at \textbf{$p$=0.08} --- i.e.\ the headline age difference is not itself
statistically resolved, and the early-onset interval very nearly touches zero. Read the
other direction instead, where the contrast is firmer: lifestyle differences map to
\textbf{3.35-fold (1.76--9.79)} larger early-onset incidence differences under high
pollution against \textbf{1.52-fold (1.33--1.74)} under low. A composite pollution score
cannot attribute anything to particles specifically.}
{10.1093/eurjpc/zwag456}

\paperentry{Pan SG et al.\ 2026}
{Ozone, not PM, survives two-pollutant adjustment against cardiac arrest --- in one city-year}
{Front Public Health}
{\textbf{7{,}996} out-of-hospital cardiac arrests from the Tianjin prehospital emergency
system, July 2023--June 2024 (mean 21.9/day), with a quasi-Poisson DLNM over lag 0--14 for
temperature and lag0--1 single- and two-pollutant models for PM$_{2.5}$, PM$_{10}$,
NO$_2$, SO$_2$, CO and O$_3$. \textbf{Ozone was the only pollutant consistently
associated}: per-IQR lag0--1, \textbf{+15.3\%} overall, men \textbf{RR 1.22 (1.08--1.39)},
women \textbf{1.22 (1.04--1.43)}. The temperature result is reported with unusual honesty
--- a monotonic cumulative risk curve, minimum-mortality temperature $-$1.9\,$^\circ$C,
apparent heat peak at lag 1 (RR $\approx$1.57 at the 95th percentile) whose
\textbf{confidence intervals include unity}, and the authors say so in the abstract. The
binding limit is exposure length: \textbf{one city, twelve months}, so season and pollutant
are almost fully confounded and the null PM result is uninformative rather than negative.}
{10.3389/fpubh.2026.1871041}

\paperentry{Jiang Z et al.\ 2026}
{Multipollutant exposure and cardio-renal-metabolic multimorbidity trajectories --- metadata only}
{Atmos Environ}
{No abstract was deposited to Crossref, Europe PMC or Semantic Scholar at index time.
Design, cohort, exposure model, sample size and setting are all unknown and are recorded
as such rather than inferred from the title. It is the \textbf{32nd}
\textit{Atmospheric Environment} record to enter this corpus abstract-free since 1 August,
which makes this a chronic publisher deposit gap rather than an incident; the DOI is
queued for the weekly re-check, and the record occupies a
\textit{Cardiovascular \& metabolic} slot on the title's face only.}
{10.1016/j.atmosenv.2026.122340}

\band{Reproductive \& developmental\hfill 2 records}{Sky}

\paperentry{Mbishi et al.\ 2026}
{Perinatal PM and infant head circumference: unambiguous, and small}
{Environ Pollut}
{\textbf{119{,}110} Flemish children with BMI-, length- and head-circumference-for-age
$z$-scores from routine well-baby clinic records, perinatal PM$_{10}$ and PM$_{2.5}$ from
the RIO detrended-kriging model, and mixed-effects models carrying a 5-df natural cubic
spline in child age with exposure-by-age interaction and child-level random intercepts and
slopes. At 12 months, per IQR: head circumference \textbf{$-$0.0499 (95\% CI $-$0.0606,
$-$0.0393)} $z$ for PM$_{10}$ and \textbf{$-$0.0645 ($-$0.0771, $-$0.0518)} for PM$_{2.5}$,
with inverse associations sustained between 6 and 24 months. BMI- and length-for-age
associations \emph{change sign} across age, and green-space associations are inconsistent
across measures. The authors state that estimates are small and that most nominally
significant associations are of limited magnitude --- roughly a fifteenth of a standard
deviation. Exposure is municipality-scale interpolation, which at this $n$ makes precision
a poor guide to accuracy.}
{10.1016/j.envpol.2026.129064}

\paperentry{Wei Q et al.\ 2026}
{Weekly component windows: black carbon at 12--15 weeks, nitrate at 5--12 and 22--24}
{Int J Epidemiol}
{\textbf{2{,}680} mother--newborn pairs and \textbf{9{,}339} growth ultrasounds (abdominal
circumference, head circumference, femur length, estimated fetal weight), Guangzhou
2017--2020, with \emph{weekly} black carbon, organic matter, nitrate, sulfate and ammonium
from the TAP model entered into distributed lag models with GEE, plus machine-learning
mixture analysis. Femur-length reduction is most strongly associated with \textbf{black
carbon at gestational weeks 12--15}; head-circumference deficits with \textbf{nitrate at
weeks 5--12 and 22--24}; estimated fetal weight with \textbf{nitrate at weeks 6--8}. Black
carbon and nitrate dominate the mixture. This is the exposure resolution the Flemish cohort
above lacks, at 45$\times$ fewer children. The threat to inference is that five components
drawn from one gridded product co-vary strongly, so ``which component'' is partly a
question about the exposure model rather than about biology; the abstract reports no effect
sizes for the identified windows.}
{10.1093/ije/dyag185}

\band{Mechanistic toxicology\hfill 2 records}{Violet}

\paperentry{Heo et al.\ 2026}
{Three weeks of PM$_{10}$ at the air--liquid interface: barrier down, differentiation lost,
Notch up}
{Tuberc Respir Dis (Seoul)}
{Differentiated human small airway epithelial ALI cultures exposed to PM$_{10}$ for 7, 14
and 21 days. Transepithelial electrical resistance \textbf{fell progressively}; club-cell
and ciliated-cell markers fell while basal-cell markers were preserved, i.e.\ a
differentiation block rather than uniform cytotoxicity. Surfactant \textbf{SP-A and SP-B
mRNA fell}, yet \textbf{secreted SP-A rose while SP-B secretion fell} --- a dissociation
the paper reports without over-explaining. IL-1$\beta$, IL-6, IL-8 and TNF-$\alpha$ rose;
cleaved Notch1 and Notch3, phospho-NF-$\kappa$B, I$\kappa$B$\alpha$, HES1 and HEY1 all
increased. The chronic, repeated-dose ALI design is a genuine advance on single-dose
submerged culture. Limits are the usual ones and are not minor: \textbf{PM$_{10}$ dose,
source and donor number are not stated in the abstract}, and no pathway inhibition
experiment establishes that Notch is causal rather than concurrent.}
{10.4046/trd.2026.0072}

\paperentry{Ren Y et al.\ 2026}
{Lung-on-chip for PM, appraised model by model}
{Research (Wash D C)}
{Critical review of \emph{in vitro} platforms for PM--lung interaction, from 2D monolayers
through 3D bioprinted constructs and organoids to lung-on-chip devices, with each class
assessed for what it can and cannot reproduce. The central argument is that
breath-mimicking cyclic mechanical strain is not an incidental refinement but a determinant
of how particles interact with the epithelial interface --- deposition, clearance and
mechanotransduction all change under motion, so static models systematically misrepresent
the dose actually delivered. Useful as an orientation to which platform answers which
question. It is a narrative review with no systematic search, no pooled comparison and a
closing pitch for the authors' own ``highly integrated and intelligent'' platform, so the
appraisal is not disinterested.}
{10.34133/research.1414}

\band{Neuro / mental health\hfill 1 record}{Deep}

\paperentry{Ren X et al.\ 2026}
{Black carbon dominates a PM$_{2.5}$ mixture on adolescent self-injury --- in cross-section}
{Front Public Health}
{\textbf{18{,}398} adolescents in grades 7--12 across 15 schools in four Chinese provinces,
with non-suicidal self-injury by ANSAQ and multidimensional biological rhythm by SBRDA.
PM$_{2.5}$ and its constituents were positively associated with NSSI risk, black carbon
most strongly: \textbf{OR 1.76 (1.60--1.94)}, and dominant in the mixture models
(\textbf{WQS weight 0.810}, Qgcomp effect 0.671). This is the largest effect on today's
forest panel and rests on the weakest design there. A single cross-sectional survey round
cannot order exposure and outcome in time; exposure is assigned by school or residence, so
the black-carbon signal is also a traffic-proximity and urbanicity signal; and rhythm
disorder is plausibly a mediator, a confounder or a consequence, which stratification
cannot distinguish. The mixture weight is the informative part --- it says the components
are not interchangeable --- and the odds ratio is not the part to quote.}
{10.3389/fpubh.2026.1887542}

\band{Respiratory \& allergic\hfill 1 record}{Sage}

\paperentry{Wangupadcha et al.\ 2026}
{A null result, reported as null, at 42\,$\mu$g/m$^3$}
{J Glob Health}
{Time-series of \textbf{924} respiratory visits (ICD-10 J00--J99, mean 6.37/day) among
residents aged \textbf{$\geq$60} living within 5\,km of a newly established industrial area
in central Thailand, November 2023--March 2024, against daily PM$_{2.5}$ from the nearest
government monitor, with negative binomial single-lag and moving-average models adjusted
for temperature, humidity, wind speed, day of week and trend. Mean PM$_{2.5}$ was
\textbf{42.43\,$\mu$g/m$^3$ (SD 15.12)}, far above WHO guidance. The largest estimate was at
lag 4: \textbf{IRR 1.055 (0.995--1.119)} per 10\,$\mu$g/m$^3$ --- positive, not significant,
and described that way by the authors rather than spun. Subgroup estimates were numerically
larger in women and at 60--69 with non-significant interaction tests. The study is
underpowered by construction: \textbf{five months, 924 events, one monitor} for a 5\,km
radius. Worth carrying as a well-behaved null and as a monitoring-density argument.}
{10.7189/jogh.16.04269}

\band{Other clinical endpoints\hfill 3 records}{Slate}

\paperentry{Ding E et al.\ 2026}
{Compound climate extremes and influenza, with a submultiplicative interaction}
{Environ Sci Technol}
{\textbf{630{,}120} laboratory-confirmed influenza cases from hospitals in \textbf{335
Chinese cities}, 2005--2019, analysed by case-crossover with distributed lag non-linear
models, machine-learning variable ranking, and an LSTM forecasting proof of concept for
Beijing. Temperature, humidity and particulate matter rank as the primary contributors.
The result worth keeping: co-exposure to multiple extreme events carries higher joint risk
than single exposure, but the interaction is \textbf{submultiplicative on the
multiplicative scale with no evidence of additive interaction} --- so the joint effect is
less than the product and the naive stacking of extreme-event risks overstates it. Adding
extreme-event terms improved daily forecasting in the Beijing model. Caveats: the abstract
reports no effect sizes at all, hospital-based case ascertainment varies across 335 cities
and 15 years, and one-city forecasting is a demonstration, not a validation.}
{10.1021/acs.est.6c04755}

\paperentry{Zheng M et al.\ 2026}
{Occupational particulate exposure ranks with alcohol in an ecological cancer model}
{Ann Med}
{GBD 2021 incidence for adults \textbf{$\geq$55} across \textbf{203 countries}, 1990--2021,
classified into quartile-based co-burden categories for oesophageal and lip/oral-cavity
cancer, with 58 candidate exposures screened by random forest and SHAP to 36, then to 25
after variance-inflation screening, entering negative binomial models with population
offsets. Countries split into consistent (73), oesophageal-dominant (61) and LOC-dominant
(69). \textbf{Occupational exposure to particulate matter, gases and fumes carries
RR 1.29 (1.13--1.49)} for oesophageal cancer, comparable to high alcohol use
(1.32, 1.13--1.53) and above low nut/seed intake (1.16). The authors state the correct
caveat themselves: these are \textbf{country-level ecological associations} and establish
no individual-level causality. GBD exposure estimates for occupational PM are modelled,
not measured, so the estimate is a model-on-model contrast.}
{10.1080/07853890.2026.2721134}

\paperentry{Wang C et al.\ 2026}
{Green space and cardiovascular health where the air-quality pathway is 75\% of the model
and the effects are fragile}
{Environ Sci Technol}
{Two-stage gradient-boosting plus structural equation modelling on a clinical cohort of
\textbf{23{,}312} patients in Nanjing. \textbf{Air-quality-related pathways account for
more than 75\% of the relative statistical contribution to model variance}, exceeding
psychological-restoration proxies --- which, if it holds, says green-space cardiovascular
benefit in a dense polluted city is mostly an air-quality mechanism. Distinct networks
emerge for chronic ischaemic heart disease (vegetation vitality inverse to NO$_2$ and
SO$_2$, linked to metabolic markers) and heart failure. The paper is unusually candid about
its own weaknesses and they are disqualifying for the headline: the macroscale aerodynamic
interpretation is called speculative for want of morphological or fluid-dynamics data, and
the age-stratified canopy-density $\times$ PM$_{2.5}$ slopes are \textbf{nonsignificant with
extremely small effect sizes}. Relative variance contribution is not effect size.}
{10.1021/acs.est.6c08296}

\vspace{4mm}

% ==================================================================== TABLE
\section*{\sffamily\bfseries\color{Deep}Corpus register}
\vspace{-1.5mm}{\color{Amber}\rule{\linewidth}{1.1pt}}\vspace{2.5mm}

{\fontsize{6.7}{8.4}\selectfont
\setlength{\tabcolsep}{2.4pt}
\begin{longtable}{@{}>{\RaggedRight}p{25mm}>{\RaggedRight}p{41mm}>{\RaggedRight}p{27mm}%
>{\RaggedRight}p{23mm}>{\RaggedRight}p{18mm}>{\RaggedRight}p{15mm}@{}}
\toprule
\textsf{\textbf{First author}} & \textsf{\textbf{Focus}} & \textsf{\textbf{Journal}} &
\textsf{\textbf{Design}} & \textsf{\textbf{Metric}} & \textsf{\textbf{Region}}\\
\midrule
\endfirsthead
\toprule
\textsf{\textbf{First author}} & \textsf{\textbf{Focus}} & \textsf{\textbf{Journal}} &
\textsf{\textbf{Design}} & \textsf{\textbf{Metric}} & \textsf{\textbf{Region}}\\
\midrule
\endhead
\multicolumn{6}{@{}l@{}}{\cellcolor{Mist}\textsf{\textbf{\textcolor{Deep}{Neuro / mental health}}}}\\[0.6mm]
\href{https://doi.org/10.1001/jamanetworkopen.2026.25015}{Dales} & Wildfire-sourced PM$_{2.5}$ and ED visits for migraine and other primary headache & \textit{JAMA Netw Open} & Case-crossover & PM$_{2.5}$ (wildfire) & Canada\\
\href{https://doi.org/10.1093/ageing/afag215}{Glover} & Umbrella meta-review of urban environmental exposures and cognitive health & \textit{Age Ageing} & Umbrella review / meta-analysis & PM$_{2.5}$ & Global\\
\href{https://doi.org/10.1016/j.nsa.2026.107019}{Goossens} & In utero air pollution exposure and pre-/postnatal brain development & \textit{Neurosci Appl} & Narrative review & PM$_{2.5}$ / PM$_{10}$ / PM$_{\mathrm{coarse}}$ & Global\\
\href{https://doi.org/10.1016/j.isci.2026.116617}{Ma} & PM$_{2.5}$-induced neurotoxicity and depression: exposure risks and mechanisms & \textit{iScience} & Narrative review & PM$_{2.5}$ + components & Global\\
\href{https://doi.org/10.3760/cma.j.cn112150-20250804-00752}{Pan} & Maternal PM exposure and offspring neurodevelopment: effects and mechanisms & \textit{Zhonghua Yu Fang Yi Xue Za Zhi} & Narrative review & PM (bulk) & China\\
\addlinespace[0.9mm]
\multicolumn{6}{@{}l@{}}{\cellcolor{Mist}\textsf{\textbf{\textcolor{Teal}{Cardiovascular \& metabolic}}}}\\[0.6mm]
\href{https://doi.org/10.3760/cma.j.cn112338-20251120-00843}{Huang} & PM$_{2.5}$ and dyslipidaemia in rural Chinese children; dietary effect modification & \textit{Zhonghua Liu Xing Bing Xue Za Zhi} & Cross-sectional (multistage) & PM$_{2.5}$ & China\\
\href{https://doi.org/10.32394/pe/218966}{Kazmierska} & Air pollution and cardiovascular disease: review of current data & \textit{Przegl Epidemiol} & Narrative review & PM$_{2.5}$ / PM$_{10}$ & Global\\
\href{https://doi.org/10.1007/s00125-026-06814-2}{Xi} & Air pollution, built environment and temperature vs. prediabetes phenotypes & \textit{Diabetologia} & Prospective cohort & PM$_{2.5}$ & Germany\\
\href{https://doi.org/10.1161/JAHA.125.046917}{Yang} & Short-term PM$_{2.5}$ and hospitalisation/ED use in atrial fibrillation & \textit{J Am Heart Assoc} & Case-crossover & PM$_{2.5}$ (wildfire-influenced) & USA\\
\addlinespace[0.9mm]
\multicolumn{6}{@{}l@{}}{\cellcolor{Mist}\textsf{\textbf{\textcolor{Sky}{Reproductive \& developmental}}}}\\[0.6mm]
\href{https://doi.org/10.1016/j.jes.2025.10.043}{Chen} & Prenatal air pollution, fetal growth trajectory phenotypes and child neurodevelopment & \textit{J Environ Sci (China)} & Prospective cohort & PM$_{2.5}$ / PM$_{10}$ & China\\
\href{https://doi.org/10.1016/j.ecoenv.2026.120566}{Huang} & Joint prenatal bisphenol + air-pollutant mixture exposure and preschool IQ & \textit{Ecotoxicol Environ Saf} & Prospective cohort (mixtures) & PM$_{2.5}$ / PM$_{10}$ & China\\
\href{https://doi.org/10.1007/s40572-026-00552-8}{Peterson} & Ambient air pollution, PAHs and pubertal development: critical review & \textit{Curr Environ Health Rep} & Critical review & PM$_{2.5}$ components & Global\\
\href{https://doi.org/10.1016/j.envres.2026.125313}{Xie} & Long-term PM$_{2.5}$ and timing of pubertal development in boys and girls & \textit{Environ Res} & Prospective cohort & PM$_{2.5}$ & China\\
\addlinespace[0.9mm]
\multicolumn{6}{@{}l@{}}{\cellcolor{Mist}\textsf{\textbf{\textcolor{Sage}{Respiratory \& allergic}}}}\\[0.6mm]
\href{https://doi.org/10.1007/s00484-026-03289-z}{Bayar Muluk} & Climate change and otolaryngology: airway disease impacts and emission reduction & \textit{Int J Biometeorol} & Narrative review & PM$_{2.5}$ & Global\\
\href{https://doi.org/10.4103/mgr.MEDGASRES-D-25-00164}{Liao} & Short-term air pollution, respiratory symptoms and lung function in asthma & \textit{Med Gas Res} & Panel study & PM$_{2.5}$ / PM$_{10}$ & China\\
\href{https://doi.org/10.3390/epidemiologia7040097}{Sanchez de Prada} & Ambient PM and asthma-coded admissions in children 0-3 y, Spain & \textit{Epidemiologia} & Ecological panel & PM$_{2.5}$ / PM$_{10}$ & Spain\\
\href{https://doi.org/10.1016/j.jhazmat.2026.143076}{Yang} & Allergenic vs. toxic risk of airborne seed hairs compared with ambient PM$_{2.5}$ & \textit{J Hazard Mater} & Field measurement & PM$_{2.5}$ + bioaerosol & China\\
\addlinespace[0.9mm]
\multicolumn{6}{@{}l@{}}{\cellcolor{Mist}\textsf{\textbf{\textcolor{Violet}{Mechanistic toxicology}}}}\\[0.6mm]
\href{https://doi.org/10.1016/j.jes.2025.12.074}{Ding} & PM$_{2.5}$ promotes vascular calcification via miR-27b-5p/HES1 signalling & \textit{J Environ Sci (China)} & In vitro + organoid + murine & PM$_{2.5}$ & China\\
\href{https://doi.org/10.1007/s00360-026-01698-5}{Maraschi} & Osmoregulatory disruption and metal bioaccumulation from settleable atmospheric PM & \textit{J Comp Physiol B} & Controlled exposure (ecotox) & Settleable APM & Brazil\\
\href{https://doi.org/10.1016/j.jes.2025.11.018}{Wu} & Macrophage NLRP3 inflammasome in real-ambient PM-induced hepatic glucose dysregulation & \textit{J Environ Sci (China)} & Real-ambient murine + co-culture & Real-ambient PM & China\\
\addlinespace[0.9mm]
\multicolumn{6}{@{}l@{}}{\cellcolor{Mist}\textsf{\textbf{\textcolor{Clay}{Occupational \& indoor}}}}\\[0.6mm]
\href{https://doi.org/10.3389/fmedt.2026.1796828}{Chen} & Nested suction instrument to cut intraoperative surgical-smoke PM$_{2.5}$ & \textit{Front Med Technol} & Bench / device experiment & PM$_{2.5}$ (surgical smoke) & China\\
\href{https://doi.org/10.1093/annweh/wxag055}{Cole} & Determinants of respirable crystalline silica exposure in Australian tunnelling & \textit{Ann Work Expo Health} & Observational exposure study & Respirable dust / RCS & Australia\\
\href{https://doi.org/10.1016/j.jhazmat.2026.143081}{Ding} & Stage-resolved hair Zn-isotope dynamics as a firefighting combustion-exposure tracer & \textit{J Hazard Mater} & Repeated-measures biomarker & Combustion-derived PM & China\\
\href{https://doi.org/10.3390/jox16040119}{Ramsden} & Xenobiotic hazards in aircraft cabin air & \textit{J Xenobiot} & Review + risk economics & Ultrafine particles & Global\\
\addlinespace[0.9mm]
\multicolumn{6}{@{}l@{}}{\cellcolor{Mist}\textsf{\textbf{\textcolor{Amber}{Sensing, forecasting \& instrumentation}}}}\\[0.6mm]
\href{https://doi.org/10.1016/j.jenvman.2026.130515}{Shi} & Adaptive spatiotemporal graph transformer for multi-site multi-step PM$_{2.5}$ forecasting & \textit{J Environ Manage} & Deep-learning model & PM$_{2.5}$ & China / India\\
\href{https://doi.org/10.3390/bios16070380}{Tsuchiya} & Seasonal and indoor-outdoor characteristics of fluorescent aerosol particles in offices & \textit{Biosensors} & Real-time sensor campaign & Bioaerosol / coarse & Japan\\
\addlinespace[0.9mm]
\multicolumn{6}{@{}l@{}}{\cellcolor{Mist}\textsf{\textbf{\textcolor{Amber}{Exposure assessment \& modelling}}}}\\[0.6mm]
\href{https://doi.org/10.1029/2025GH001621}{Solorzano-Araque} & Multiscale exposure modelling: LOTOS-EUROS + mobility surveys + particle toxicity & \textit{GeoHealth} & CTM + mobility integration & PM$_{2.5}$ / PM$_{10}$ & Colombia\\
\addlinespace[0.9mm]
\multicolumn{6}{@{}l@{}}{\cellcolor{Mist}\textsf{\textbf{\textcolor{Coral}{Burden, policy \& mitigation}}}}\\[0.6mm]
\href{https://doi.org/10.1002/hsr2.72886}{Karamzad} & Ischaemic heart disease burden in MENA 1990-2021; PM attributable fraction 28.2\% & \textit{Health Sci Rep} & GBD secondary analysis & Ambient PM & MENA\\
\href{https://doi.org/10.1371/journal.pone.0347989}{Rees} & HEPA filtration to reduce respiratory infection in care homes & \textit{PLoS One} & RCT process evaluation & Filtered PM & UK\\
\href{https://doi.org/10.1016/j.jes.2026.01.009}{Shan} & Policy-driven emission reductions and post-pandemic rebound, high-resolution inventories & \textit{J Environ Sci (China)} & Emission inventory & PM emissions & China\\
\addlinespace[0.9mm]
\multicolumn{6}{@{}l@{}}{\cellcolor{Mist}\textsf{\textbf{\textcolor{Slate}{Other clinical endpoints}}}}\\[0.6mm]
\href{https://doi.org/10.3390/dj14070440}{Lupita} & PM$_{2.5}$, preventive-care capacity and caries/edentulism burden across the EU & \textit{Dent J} & Ecological panel & PM$_{2.5}$ & EU-27\\
\href{https://doi.org/10.1016/j.jdin.2026.06.008}{Park} & Long-term PM exposure and risk of androgenetic alopecia & \textit{JAAD Int} & Nationwide cohort & PM$_{2.5}$ / PM$_{10}$ & South Korea\\
\href{https://doi.org/10.3171/2026.3.PEDS2615}{Sun} & Environmental exposure and area deprivation vs. paediatric hydrocephalus outcomes & \textit{J Neurosurg Pediatr} & Registry cohort + geospatial & PM$_{2.5}$ (satellite-derived) & USA\\
\addlinespace[0.9mm]
\bottomrule
\end{longtable}}


\vspace{4mm}
\begin{keybox}
{\sffamily\bfseries\footnotesize\color{Deep}Provenance and method}\\[1.2mm]
{\fontsize{7.2}{8.8}\selectfont
Entry window \textbf{4 September 2026}, single day, opened against a
\texttt{last\_entry\_date} of 3 September. Sources queried on 5 September 2026 at 11:4x
CDT; under this pipeline's 22:00-local rule, 4 September is the newest buildable date and
5 September is not opened. Exactly one day was owed and no gap remains. \textbf{The W36
weekly rollup (30 August--5 September) is therefore not built here}: it needs 5 September,
which only the 23:00 run can open.
\textbf{PubMed} E-utilities on \texttt{[EDAT]} 2026/09/04: \textbf{14} hits on the health
axis, \textbf{3} on the instrumentation axis; the PubMed connector, run independently with
looser field tags, returned 15 and 4 and surfaced no record the harvester had missed.
\textbf{Europe PMC} \texttt{CREATION\_DATE}: 24.
\textbf{Crossref} by-ISSN, windowed on \texttt{created}: \textbf{57} deposits.
\textbf{arXiv} answered without an HTTP 429 for the second consecutive run (149\,kB of
Atom) but its newest submission remains 2026-08-20, so the undated relevance query again
contributed nothing; \textbf{OpenAlex} returned HTTP 429, the plan-gated date filter having
been broken since 28 July. \textbf{Consensus} returned three low-cost-sensor calibration
papers and all three were checked against Crossref \texttt{created} rather than assumed:
two are out of window (2026-03-26, 2026-04-21) and the third is already in this corpus from
an earlier issue --- a sixth consecutive nil, now verified rather than inspected.
After collapsing cross-source duplicates by PMID, DOI and normalised title, \textbf{91
unique records} were considered; \textbf{9 were already shipped} in earlier issues,
\textbf{27 were screened in scope} and \textbf{55 were screened out}. Notable exclusions:
two \textit{International Journal of Population Data Science} conference abstracts, on
London Underground commuting and on ambient air pollution and dementia --- both squarely on
topic, both with no DOI and no retrievable text; two \textit{medRxiv} preprints on thermal
variability and childhood respiratory mortality, where PM$_{2.5}$ enters only as an
adjustment covariate; an Apple Hearing Study analysis of noise and heart rate variability,
methodologically close to distributed personal sensing but carrying no air-pollution
metric; and a critical review of black-carbon marine carbon sequestration, which is
atmospheric in origin and oceanographic in question.
\textbf{Trial watch:} one new interventional registration, \textbf{NCT07805824} (cluster
RCT, workplace heat and air-pollution mitigation, 20 restaurants and 200 oven workers,
Najaf, Iraq; primary endpoint an objective gingival melanin index, sole secondary endpoint
change in personal PM$_{2.5}$ by gravimetric and real-time monitoring).
\texttt{analyze\_endpoints} was run because this is a new registration and not an update.
Its phase label, start and completion dates are internally inconsistent and are recorded as
unreliable. Logged to \texttt{state/trials.json} for the W36 weekly.
The DOI gate passed all 27 records with \textbf{0 failures and 0 warnings}. All 24 distinct
design labels were already present in the figure map, so no \texttt{plots.py} edit was made
this run.
Subtopic, design, particle-metric and geography labels are single-label by dominant
contribution and are a judgement, not a controlled vocabulary. All DOI links resolve to the
publisher of record; abstracts remain the copyright of their respective publishers.\par}
\end{keybox}

\end{document}
